Official marking scheme:

\begin{description}
\item[(1)] Dec and RA grids made: full score.
\item[(2)] Smallest circle, full score; medium circle, half score; largest
  circle, quarter score (concentric tolerance circles are drawn around each
  of the five stars on the key chart).
\item[(3)] Smallest circle, full score; medium circle, half score; largest
  circle, quarter score (tolerance circles around the Moon's position).
\item[(4)] Smallest circle, full score; medium circle, half score; largest
  circle, quarter score (tolerance circles around each Messier object).
\item[(5)] Galactic plane inside the yellow line: full score. Galactic
  plane crosses the yellow line: half score.
\item[(6)] $\pm 2^\circ$, full score; $\pm 5^\circ$, half score.
\end{description}

% FIGURE NEEDED: key chart with the five brightest stars circled and ranked
% 1-5 in red and the Moon's position for 29 July 2015 circled in green,
% north-east of Nova Sagittarii near the M25/M17 region
% (2015_SKYMAP_sol.pdf page 2). Read from the chart, the ranked stars are:
% 1 = lambda Sco (Shaula), 2 = epsilon Sgr (Kaus Australis),
% 3 = theta Sco (Sargas), 4 = sigma Sgr (Nunki), 5 = kappa Sco.

% FIGURE NEEDED: key chart with every Messier object in the field circled
% in red (any five accepted) (2015_SKYMAP_sol.pdf page 3)

% FIGURE NEEDED: key chart with the accepted corridor for the galactic
% plane drawn as a thick yellow band (2015_SKYMAP_sol.pdf page 5)

The key lists the Messier objects in the field of view (any five could be
marked for task 4):

\begin{center}
\begin{tabular}{ccc}
Messier & Magnitude & size ($'$) \\
\hline
7 & 3.3 & 80 \\
6 & 4.2 & 15 \\
25 & 4.6 & 32 \\
22 & 5.2 & 24 \\
23 & 5.5 & 27 \\
8 & 5.8 & 90 \\
21 & 5.9 & 13 \\
16 & 6 & 35 \\
17 & 6 & 46 \\
20 & 6.3 & 29 \\
62 & 6.6 & 14 \\
20 & 6.9 & 11 \\ % sic -- M20 is listed twice in the key table
18 & 6.9 & 9 \\
55 & 7 & 19 \\
19 & 7.2 & 13 \\
69 & 7.7 & 7 \\
54 & 7.7 & 9 \\
9 & 7.9 & 9 \\
70 & 8.1 & 7 \\
24 & 11 & 5 \\ % sic -- M24's listed magnitude and size as in the key
\end{tabular}
\end{center}
